\chapter{BKNR Import/Export}

\vbox{
    \centering
    \includegraphics{impexicon}
    \vspace{1cm}}

\section{ Introduction}

We often have to communicate data with the external world, for
example sharing data with external programs or getting information
from a third-party. The XML format has been pretty successful as a
common base format for digital information. However, parsing and
handling and serializing XML data is not very practical, though
not very difficult, and most of the time just plain
boring. Imagine you have a XML file containing information about
books in a library. Most of the time, you will parse the XML data,
and instantiate CLOS objects in order to use standard LISP
functions on the data. The parsing and serializing code is most of
the time trivial. The XML import/export module in BKNR enables you
to annotate standard CLOS class definitions in order to make
parsing and serializing XML data automatic.

The XML import/export modules works by using DTDs to represent the
format of the XML file. DTDs are quite simple and pretty
widespread. A DTD specifies which elements are valid in a XML
file, which attributes these elements have, and how the elements
are nested. When defining a CLOS class, you can specify a DTD, and
annotate the slot definitions with information about which slot is
mapped to which attribute or element in the DTD.

This tutorial will show you how to parse and serialize CLOS
classes to XML, and how to use features like specifying relations
on IDs.


\section{ Obtaining and loading BKNR XML import/export}

You can obtain the current CVS sources of BKNR by following the
instructions at `http://bknr.net/blog/bknr-devel'. Add the `src'
directory of BKNR to your `asdf:*central-registry*', and load the
indices module by evaluating the following form:

\begin{Verbatim}[fontsize=\small,frame=leftline,framerule=0.9mm,rulecolor=\color{gray},framesep=5.1mm,xleftmargin=5mm,fontfamily=cmtt]
(asdf:oos 'asdf:load-op :bknr-impex)
\end{Verbatim}
Then switch to the `bknr.impex' package to try out the tutorial.

\begin{Verbatim}[fontsize=\small,frame=leftline,framerule=0.9mm,rulecolor=\color{gray},framesep=5.1mm,xleftmargin=5mm,fontfamily=cmtt]
(in-package :bknr.impex)
\end{Verbatim}


\section{ A simple example}

We start with our book example. We have to manipulate informations
about the books present in a library, and the library provides us
with a XML file containing all the books present in the library. A
sample XML file has the following format:

\begin{Verbatim}[fontsize=\small,frame=leftline,framerule=0.9mm,rulecolor=\color{gray},framesep=5.1mm,xleftmargin=5mm,fontfamily=cmtt]
<books>
  <book id="1" isbn="234114">
    <author>J.R.R Tolkien</author>
    <title>The Lord of the Rings</title>
  </book>
  <book id="2" isbn="235431">
    <author>William Gibson</author>
    <title>Neuromancer</title>
  </book>
</books>
\end{Verbatim}
The DTD for the sample above is pretty simple too:

\begin{Verbatim}[fontsize=\small,frame=leftline,framerule=0.9mm,rulecolor=\color{gray},framesep=5.1mm,xleftmargin=5mm,fontfamily=cmtt]
<!ELEMENT books (book)*>
<!ELEMENT book (author,title)>
<!ELEMENT author (#PCDATA)>
<!ELEMENT title (#PCDATA)>

<!ATTLIST book id   ID    #REQUIRED
               isbn CDATA #REQUIRED>
\end{Verbatim}
We now specify our CLOS class book without XML annotations. The
common approach to XML parsing would now implement either a SAX
parser filling a `BOOK' class from a `book' XML element, or a DOM
parser walking down the DOM tree to instantiate `BOOK' objects at
`book' leaves.

\begin{Verbatim}[fontsize=\small,frame=leftline,framerule=0.9mm,rulecolor=\color{gray},framesep=5.1mm,xleftmargin=5mm,fontfamily=cmtt]
(defclass book ()
  ((author :initarg :author :reader book-author)
   (id :initarg :id :reader book-id :type integer)
   (isbn :initarg :isbn :reader book-isbn)
   (title :initarg :title :reader book-title)))

(defmethod print-object ((book book) stream)
  (print-unreadable-object (book stream :type t :identity t)
    (format stream "~S" (book-title book))))
\end{Verbatim}
We can however extend the `BOOK' class definitions by giving it
the metaclass `XML-CLASS', by specifying the DTD to be used by the
class, and by specifying the XML element corresponding to the
class. We also annotate the slot definitions.

\begin{Verbatim}[fontsize=\small,frame=leftline,framerule=0.9mm,rulecolor=\color{gray},framesep=5.1mm,xleftmargin=5mm,fontfamily=cmtt]
(defvar *tutorial-dtd*
  (cxml:parse-dtd-file "xml-impex/tutorial.dtd"))

(defclass book ()
  ((author :initarg :author :reader book-author
      :element "author")
   (id :initarg :id :reader book-id :type integer
       :attribute "id" :parser #'parse-integer)
   (isbn :initarg :isbn :reader book-isbn
    :attribute "isbn")
   (title :initarg :title :reader book-title
     :element "title"))
  (:metaclass xml-class)
  (:dtd *tutorial-dtd*)
  (:element "book"))
\end{Verbatim}
We can now read the XML file containing the book definitions. As
we haven't specified a CLOS class for the root element `books', we
get back a list containing all the children nodes of `books' under
the keyword `:BOOK'.

\begin{Verbatim}[fontsize=\small,frame=leftline,framerule=0.9mm,rulecolor=\color{gray},framesep=5.1mm,xleftmargin=5mm,fontfamily=cmtt]
(parse-xml-file "xml-impex/tutorial.xml" (list (find-class 'book)))
; => (:BOOK
;     (#<BOOK ("The Lord of the Rings") {4922BCD5}>
;      #<BOOK ("Neuromancer") {4922E0B5}>))
(setf *books*
      (getf (parse-xml-file "xml-impex/tutorial.xml"
             (list (find-class 'book)))
       :BOOK))
; => (#<BOOK ("The Lord of the Rings") {4922BCD5}>
;     #<BOOK ("Neuromancer") {4922E0B5}>)
\end{Verbatim}
In the same way, we can serialize our books back to XML using the
function `WRITE-TO-XML'. As we still do not have a class for the
root element (and we never will), we can specify the name of the
root element as an argument.

\begin{Verbatim}[fontsize=\small,frame=leftline,framerule=0.9mm,rulecolor=\color{gray},framesep=5.1mm,xleftmargin=5mm,fontfamily=cmtt]
(write-to-xml *books* :name "books")
; => <books>
;     <book id="1" isbn="234114">
;        <author>J.R.R Tolkien</author>
;        <title>The Lord of the Rings</title>
;     </book>
;     <book id="2" isbn="235431">
;        <author>William Gibson</author>
;        <title>Neuromancer</title>
;     </book>
;    </books>
\end{Verbatim}
As you have noticed, we use an in ID in the book class. It would
be nice to have an index for all the books. This is very simple,
as a class with metaclass `XML-CLASS' also support BKNR
indices. Thus, you can write:

\begin{Verbatim}[fontsize=\small,frame=leftline,framerule=0.9mm,rulecolor=\color{gray},framesep=5.1mm,xleftmargin=5mm,fontfamily=cmtt]
(defclass book ()
  ((author :initarg :author :reader book-author
      :element "author"
      :index-type keyword-index :index-initargs (:test #'equal)
      :index-reader books-with-author
      :index-keys all-authors)
   (id :initarg :id :reader book-id :type integer
       :attribute "id" :parser #'parse-integer
       :index-type slot-index :index-reader book-with-id
       :index-values all-books)
   (isbn :initarg :isbn :reader book-isbn
    :attribute "isbn"
    :index-type slot-index :index-initargs (:test #'equal)
    :index-reader book-with-isbn)
   (title :initarg :title :reader book-title
     :element "title"))
  (:metaclass xml-class)
  (:dtd *tutorial-dtd*)
  (:element "book"))
\end{Verbatim}
We can now import our XML file and the indices will automatically
get filled. The indices change nothing to the serialization.

\begin{Verbatim}[fontsize=\small,frame=leftline,framerule=0.9mm,rulecolor=\color{gray},framesep=5.1mm,xleftmargin=5mm,fontfamily=cmtt]
(parse-xml-file "xml-impex/tutorial.xml" (list (find-class 'book)))
; => (:BOOK
;     (#<BOOK "The Lord of the Rings" {49224D25}>
;      #<BOOK "Neuromancer" {492272CD}>))
(all-authors)
; => ("J.R.R Tolkien" "William Gibson")
(all-books)
; => (#<BOOK "The Lord of the Rings" {49224D25}>
;     #<BOOK "Neuromancer" {492272CD}>)
(books-with-author (first (all-authors)))
; => (#<BOOK "The Lord of the Rings" {49224D25}>)
; T
(book-with-id 1)
; => #<BOOK "The Lord of the Rings" {49224D25}>
;  T
(write-to-xml *books* :name "books")
; => <books>
;       <book id="1" isbn="234114">
;          <author>J.R.R Tolkien</author>
;          <title>The Lord of the Rings</title>
;       </book>
;       <book id="2" isbn="235431">
;          <author>William Gibson</author>
;          <title>Neuromancer</title>
;       </book>
;    </books>
\end{Verbatim}
Note that you have to clear the indices if you reload the same XML
file multiple times.

\begin{Verbatim}[fontsize=\small,frame=leftline,framerule=0.9mm,rulecolor=\color{gray},framesep=5.1mm,xleftmargin=5mm,fontfamily=cmtt]
(clear-class-indices (find-class 'book))
; => NIL
\end{Verbatim}


\section{ Specifying a mapping from CLOS to XML}

In the following text, ``XML class'' refers to a CLOS class with
metaclass `XML-CLASS', and ``XML object'' refers to an instance of
an ``XML class''.



\subsection{ Class options}

Specifying `XML-CLASS' as a metaclass for a CLOS class has two
effects. First, you can use slot indices and class indices for
your class, as `XML-CLASS' inherits from `INDEXED-CLASS' (see the
BKNR Indices tutorial for more information). Second, you can
specify a DTD to use with your class, and an optional element name
to which the class will be mapped. Please note that when
you have multiple XML classes, it is best to specify the same DTD
object, as `EQL' is used to reduce the list of DTDs when exporting
XML objects.

A CLOS class is mapped to an element of the DTD. In our example
above, the class `BOOK' was mapped to the DTD element `book'. This
is done by specifying the `:ELEMENT' class option. The XML
import/export code then searchs the element definition in the DTD,
and stores it in the class object. It then maps the slots of the
class to the attributes and child element of the element
definition in the DTD.



\subsubsection{ Specifying an empty class element}

The `:ELEMENT' class option is optional. The optional element
specification is there to be able to derive objects which have the
same attributes without having to rewrite everything. For example,
suppose we have the following DTD:

\begin{Verbatim}[fontsize=\small,frame=leftline,framerule=0.9mm,rulecolor=\color{gray},framesep=5.1mm,xleftmargin=5mm,fontfamily=cmtt]
<!ELEMENT test EMPTY>
<!ELEMENT test2 EMPTY>
<!ELEMENT test3 EMPTY>

<!ATTLIST test id ID #REQUIRED>
<!ATTLIST test2 id ID #REQUIRED>
<!ATTLIST test3 id ID #REQUIRED>
\end{Verbatim}
We can then write the following class definitions:

\begin{Verbatim}[fontsize=\small,frame=leftline,framerule=0.9mm,rulecolor=\color{gray},framesep=5.1mm,xleftmargin=5mm,fontfamily=cmtt]
(defvar *test-dtd*
  (cxml:parse-dtd-file #p"xml-impex/tutorial2.dtd"))

(defclass test-object ()
  ((id :initarg :id :attribute "id"
       :parser #'parse-integer
       :index-type slot-index :index-reader object-with-id
       :index-values all-objects))
  (:metaclass xml-class)
  (:dtd *test-dtd*)
  (:element nil))

(defmethod print-object ((object test-object) stream)
  (print-unreadable-object (object stream :type t)
    (format stream "~A" (slot-value object 'id))))

(defclass test (test-object)
  ((id :index-type slot-index
       :index-reader test-with-id
       :index-values all-tests))
  (:metaclass xml-class)
  (:dtd *test-dtd*)
  (:element "test"))

(defclass test2 (test-object)
  ((id :index-type slot-index
       :index-reader test2-with-id
       :index-values all-test2s))
  (:metaclass xml-class)
  (:dtd *test-dtd*)
  (:element "test2"))

(defclass test3 (test-object)
  ((id :index-type slot-index
       :index-reader test3-with-id
       :index-values all-test3s))
  (:metaclass xml-class)
  (:dtd *test-dtd*)
  (:element "test3"))
\end{Verbatim}
When we parse a sample file, we get the following results:

\begin{Verbatim}[fontsize=\small,frame=leftline,framerule=0.9mm,rulecolor=\color{gray},framesep=5.1mm,xleftmargin=5mm,fontfamily=cmtt]
(parse-xml-file "xml-impex/tutorial2.xml"
      (mapcar #'find-class '(test test2 test3)))
; => (:TEST3 (#<TEST3 3>) :TEST2 (#<TEST2 2>)
;     :TEST (#<TEST 1>))
(all-objects)
; => (#<TEST 1> #<TEST2 2> #<TEST3 3>)
(all-tests)
; => (#<TEST 1>)
(object-with-id 1)
; => #<TEST 1>
;    T
(test2-with-id 2)
; => #<TEST2 2>
;    T
\end{Verbatim}


\subsection{ Slot options}

The slots of a class can be mapped to different structures of the
DTD. A slot can be:



\subsubsection{ The :attribute slot option}

an ``attribute'' - the value of the slot is stored as an attribute
of the element corresponding to the class. For example, in the
book example, both the slot `ID' and the slot `ISBN' are stored as
attribute of the element ``book''. The import/export code verifies
that the attributes are specified in the DTD.



\subsubsection{ The :element slot option}

an ``element'' - the value of the slot is stored as the CDATA body
of a child element. For example, in the book example, both the
slot `TITLE' and the slot `AUTOR' are stored as child elements of
the element ``book''. The element has to be specified as a single or
optional child of the class element. When the element is specified
as multiple children (either ``*'' or ``+'' in the DTD), then the slot
value is assumed to be a list of children. For example, if we have
the DTD:

\begin{Verbatim}[fontsize=\small,frame=leftline,framerule=0.9mm,rulecolor=\color{gray},framesep=5.1mm,xleftmargin=5mm,fontfamily=cmtt]
<!ELEMENT adult (child)*>
<!ELEMENT child EMPTY>
<!ATTLIST adult name CDATA #REQUIRED>
<!ATTLIST child name CDATA #REQUIRED>
\end{Verbatim}
we can write the following class definition:

\begin{Verbatim}[fontsize=\small,frame=leftline,framerule=0.9mm,rulecolor=\color{gray},framesep=5.1mm,xleftmargin=5mm,fontfamily=cmtt]
(defvar *adult-dtd*
  (cxml:parse-dtd-file "xml-impex/tutorial3.dtd"))

(defclass adult ()
  ((name :initarg :name :attribute "name"
    :reader adult-name)
   (children :initarg :children :element "child"
        :reader adult-children))
  (:metaclass xml-class)
  (:dtd *adult-dtd*)
  (:element "adult"))

(defmethod print-object ((adult adult) stream)
  (print-unreadable-object (adult stream :type t)
    (princ (adult-name adult) stream)))

(defclass child ()
  ((name :initarg :name :attribute "name"
    :reader child-name))
  (:metaclass xml-class)
  (:dtd *adult-dtd*)
  (:element "child"))

(defmethod print-object ((child child) stream)
  (print-unreadable-object (child stream :type t)
    (princ (child-name child) stream)))
\end{Verbatim}
We can then parse the following XML file:

\begin{Verbatim}[fontsize=\small,frame=leftline,framerule=0.9mm,rulecolor=\color{gray},framesep=5.1mm,xleftmargin=5mm,fontfamily=cmtt]
<family>
  <adult name="Clara">
    <child name="Robert"/>
    <child name="Anton"/>
  </adult>
  <adult name="Joseph">
    <child name="Ludwig"/>
    <child name="Benediktine"/>
  </adult>
</family>

(setf *adults* 
      (getf (parse-xml-file "xml-impex/tutorial3.xml"
             (mapcar #'find-class '(adult child)))
       :adult))
; => (#<ADULT Clara> #<ADULT Joseph>)
(adult-children (first *))
; => (#<CHILD Robert> #<CHILD Anton>)
(write-to-xml *adults* :name "family")
; => <family>
;       <adult name="Joseph">
;          <child name="Benediktine"/>
;          <child name="Ludwig"/>
;       </adult>
;       <adult name="Clara">
;          <child name="Anton"/>
;          <child name="Robert"/>
;       </adult>
;    </family>
\end{Verbatim}


\subsubsection{ The :parent slot option}

Assume that we want to know who the parent of a child is. We can
do this by creating a slot with the `:PARENT' slot option. This
will automatically get filled by the object representing the
parent element in the XML file.

\begin{Verbatim}[fontsize=\small,frame=leftline,framerule=0.9mm,rulecolor=\color{gray},framesep=5.1mm,xleftmargin=5mm,fontfamily=cmtt]
(defclass child ()
  ((name :initarg :name :attribute "name"
    :reader child-name)
   (parent :initarg :parent :parent t
      :reader child-parent))
  (:metaclass xml-class)
  (:dtd *adult-dtd*)
  (:element "child"))

(setf *adults* 
      (getf (parse-xml-file "xml-impex/tutorial3.xml"
             (mapcar #'find-class '(adult child)))
       :adult))
; => (#<ADULT Joseph> #<ADULT Clara>)
(adult-children (first *adults*))
; => (#<CHILD Benediktine> #<CHILD Ludwig>)
(child-parent (first *))
; => #<ADULT Joseph>
\end{Verbatim}


\subsubsection{ The :body slot option}

Sometimes, an XML element has a important string as body, and we
would like to store this string in a slot. This can be done by
setting the `:BODY' slot option. Note that only one PCDATA body is
permitted when using this option. For example, we could use this
code to read in the DTD:

\begin{Verbatim}[fontsize=\small,frame=leftline,framerule=0.9mm,rulecolor=\color{gray},framesep=5.1mm,xleftmargin=5mm,fontfamily=cmtt]
<!ELEMENT book-resume (#PCDATA)>
<!ATTLIST book-resume id ID #REQUIRED
                      book-id CDATA #REQUIRED
                      reviewer CDATA #REQUIRED>

(defvar *resume-dtd*
  (cxml:parse-dtd-file "xml-impex/tutorial4.dtd"))

(defclass book-resume ()
  ((id :initarg :id :attribute "id"
       :parser #'parse-integer
       :reader book-resume-id)
   (book-id :initarg :book-id :attribute "book-id"
       :parser #'parse-integer
       :reader book-resume-book-id)
   (reviewer :initarg :reviewer :attribute "reviewer"
        :reader book-resume-reviewer)
   (review :initarg :review :body t
      :reader book-resume-review))
  (:metaclass xml-class)
  (:dtd *resume-dtd*)
  (:element "book-resume"))
\end{Verbatim}
Parsing the following file gives the results:

\begin{Verbatim}[fontsize=\small,frame=leftline,framerule=0.9mm,rulecolor=\color{gray},framesep=5.1mm,xleftmargin=5mm,fontfamily=cmtt]
<resumes>
  <book-resume id="1" book-id="45"
               reviewer="Henry von der Grande">
Bla bla bla. Resume highlight blablabla foobar blorg.
  </book-resume>
  <book-resume id="2" book-id="1337"
               reviewer="L0rd 3v1l">
Bl4 bl4 bl4. R3s\/m3 h1ghl1ght bl4bl4bl4 f00b4r bl0rg.
  </book-resume>
</resumes>

(setf *resumes*
      (getf (parse-xml-file "xml-impex/tutorial4.xml"
             (list (find-class 'book-resume)))
       :book-resume))
; => (#<BOOK-RESUME {4947F22D}> #<BOOK-RESUME {4947DB95}>)
(book-resume-review (first *))
; => "
; Bl4 bl4 bl4. R3s\\/m3 h1ghl1ght bl4bl4bl4 f00b4r bl0rg.
; "
(write-to-xml *resumes* :name "resumes")
; => <resumes>
;       <book-resume book-id="1337" id="2"
;                    reviewer="L0rd 3v1l">
;    Bl4 bl4 bl4. R3s\/m3 h1ghl1ght bl4bl4bl4 f00b4r bl0rg.
;       </book-resume>
;       <book-resume book-id="45" id="1"
;                    reviewer="Henry von der Grande">
;    Bla bla bla. Resume highlight blablabla foobar blorg.
;      </book-resume>
;    </resumes>
\end{Verbatim}


\subsection{ Using IDs in XML import/export}

Often, you want to reference objects by ID in the XML file. For
example, we could specify both authors and books on the toplevel,
and reference the author by ID inside the document definition. The
DTD would look like this:

\begin{Verbatim}[fontsize=\small,frame=leftline,framerule=0.9mm,rulecolor=\color{gray},framesep=5.1mm,xleftmargin=5mm,fontfamily=cmtt]
<!ELEMENT book (title)>
<!ATTLIST book id ID #REQUIRED
               autor CDATA #REQUIRED>
<!ELEMENT title #PCDATA>
<!ELEMENT author (name)>
<!ATTLIST author id ID #REQUIRED>
<!ELEMENT name #PCDATA>
\end{Verbatim}
We can write the following class definitions:

\begin{Verbatim}[fontsize=\small,frame=leftline,framerule=0.9mm,rulecolor=\color{gray},framesep=5.1mm,xleftmargin=5mm,fontfamily=cmtt]
(defparameter *book2-dtd*
  (cxml:parse-dtd-file "xml-impex/tutorial5.dtd"))

(defclass author ()
  ((id :initarg :id :reader author-id
       :attribute "id" :parser #'parse-integer
       :index-type slot-index :index-reader author-with-id
       :index-values all-authors)
   (name :initarg :name :reader author-name
    :element "name"))
  (:metaclass xml-class)
  (:dtd *book2-dtd*)
  (:element "author"))

(defmethod print-object ((author author) stream)
  (print-unreadable-object (author stream :type t)
    (princ (author-name author) stream)))

(defclass book ()
  ((id :initarg :id :reader book-id
       :attribute "id" :parser #'parse-integer
       :index-type slot-index :index-reader book-with-id
       :index-values all-books)
   (author :initarg :author :reader book-author
     :attribute "author" :parser #'parse-integer
     :id-to-object #'author-with-id
     :object-to-id #'author-id)
   (title :initarg :title :reader book-title
     :element "title"))
  (:metaclass xml-class)
  (:dtd *book2-dtd*)
  (:element "book"))
\end{Verbatim}
We can then read the following XML file:

\begin{Verbatim}[fontsize=\small,frame=leftline,framerule=0.9mm,rulecolor=\color{gray},framesep=5.1mm,xleftmargin=5mm,fontfamily=cmtt]
<library>
  <author id="1">
    <name>J.R.R Tolkien</name>
  </author>
  <author id="2">
     <name>William Gibson</name>
  </author>
  <book id="3" author="1">
    <title>Lord of the Rings</title>
  </book>
  <book id="4" author="2">
    <title>Neuromancer</title>
  </book>
</library>

(map nil #'clear-class-indices
     (mapcar #'find-class '(book author)))
; => NIL
(parse-xml-file "xml-impex/tutorial5.xml"
      (mapcar #'find-class '(book author)))
; => (:BOOK
;     (#<BOOK "Neuromancer" {494AEFC5}>
;      #<BOOK "Lord of the Rings" {494AD48D}>)
;     :AUTHOR (#<AUTHOR William Gibson>
;              #<AUTHOR J.R.R Tolkien>))
(all-authors)
; => (#<AUTHOR J.R.R Tolkien> #<AUTHOR William Gibson>)
(all-books)
; => (#<BOOK "Lord of the Rings" {494AD48D}>
;     #<BOOK "Neuromancer" {494AEFC5}>)
(book-author (first *))
; => #<AUTHOR J.R.R Tolkien>
(write-to-xml (append (all-books) (all-authors))
         :name "library")
; => <library>
;       <book author="1" id="3">
;          <title>Lord of the Rings</title>
;       </book>
;       <book author="2" id="4">
;          <title>Neuromancer</title>
;       </book>
;       <author id="1">
;          <name>J.R.R Tolkien</name>
;       </author>
;       <author id="2">
;          <name>William Gibson</name>
;       </author>
;    </library>
\end{Verbatim}
